\documentclass[14pt,usepdftitle=false,aspectratio=169]{beamer}
\usepackage{preambule}
\setbeamercolor{structure}{fg=black}
\usepackage[nopar]{bigoperators}\DeclareMathOperator{\oldcl}{Cl}\newcommand{\cl}[1]{\oldcl\l#1\r}\DeclareMathOperator{\oldnrm}{nrm}\newcommand\inrm[1]{\oldnrm\l#1\r}\DeclareMathOperator{\olddisc}{disc}\newcommand{\disc}[1]{\olddisc\l#1\r}\usepackage{complexes,galois,dsft}
\hypersetup{pdftitle={Théorie algébrique des nombres -- Fonction zêta de Dedekind}}
\title{Théorie algébrique des nombres\\\emph{Fonction zêta de Dedekind}}
\author{}
\date{}
\begin{document}
\begin{frame}
    \titlepage
\end{frame}
\slideq{$\zeta_K\l s\r$}{1/2}
\slider{$\sum{\substack{I\le\calO_K\\I\neq\set0}}{}{\frac1{\inrm I^s}}$}{1/2}
\slideq{Théorème de Hecke sur la convergence de la fonction zêta}{2/2}
\slider{La série $\zeta_K$ converge pour $\Re s>1$ et $\zeta_K$ se prolonge en une fonction méromorphe sur $\set{\Re s>1-\oldfrac1{\edeg K\bbQ}}$ avec un unique pôle simple en $s=1$, de résidu $\kappa=\oldfrac{2^{r_1}\l2\rmpi\r^{r_2}h_KR_K}{\omega_K\sqrt{\left|\disc K\right|}}$ où $r_1$ et $r_2$ sont le nombre de plongements réels et complexes conjugués, $h_K=\left|\cl{\calO_K}\right|$, $R_K$ le régulateur de $\calO_K$ et $\omega_K=\left|\mu_K\right|$}{2/2}
\end{document}